% =============================================================================
% Introduction, version 3 (26 November 2026)
% Working title: How asymmetric reputation updating sustains cooperation
%
% Drop-in replacement for the \section*{Introduction} block of
% "Manuscript_full .tex" (lines 114-133 in revision 3). Five paragraphs:
%   1 reputation and cooperation; 2 why reliable reputations are hard;
%   3 assessment versus updating; 4 asymmetric updating and the gap;
%   5 here we show.
% Citation keys resolve against the uploaded references_2.bib plus
% Introduction_V3_new_refs.bib (panchanathan2004, rankin2009). The repository
% copy of references.bib is older and lacks several keys used here
% (alexander1987, sugden1986, kandori1992, brandt2004, uchida2013, yoeli2013).
% Notes, flagged claims and a readable version:
% manuscript_planning/Introduction_V3_notes.md
% British spelling (Nature house style); the rest of the manuscript still uses
% American spelling.
% =============================================================================
\section*{Introduction}

% -- 1. Reputation links past behaviour to future cooperation
The consequences of helping or refusing to help rarely end with the people directly involved. Others observe the interaction or hear about it, form a judgement of the actor and adjust how they treat that person in later encounters\cite{alexander1987,milinski2016}. Reputation thereby links past behaviour to future opportunities, creating incentives to help that do not depend on meeting the same partner again. This mechanism, known as indirect reciprocity, can sustain cooperation among individuals who rarely or never interact twice\cite{nowak1998,nowak2005}. Theoretical work has shown that such cooperation can be stable\cite{sugden1986,kandori1992,leimar2001,panchanathan2003} and can even support contributions to collective action\cite{panchanathan2004}. Experiments show that people preferentially help those who have helped others\cite{wedekind2000}, that reputational concerns sustain contributions to public goods\cite{milinski2002} and that making behaviour observable increases cooperation in the field\cite{yoeli2013}. Reputation systems likewise sustain exchange between strangers online\cite{tadelis2016}. Yet indirect reciprocity is only as effective as the judgements on which it rests. Observers must reliably distinguish those who cooperate from those who do not, often from information that is fragmentary, indirect and occasionally wrong.

% -- 2. Norms, private assessment, and the two requirements of reliable reputations
Much of the theory of indirect reciprocity has asked how such judgements should be made. A social norm specifies how an observed action is evaluated in its context, for instance whether refusing to help someone in bad standing is justified or deserves condemnation\cite{ohtsuki2004,brandt2004}. An exhaustive search identified eight norms, the ``leading eight'', that sustain cooperation when reputations are public, even if errors occur\cite{ohtsuki2004,ohtsuki2006}. When observers instead hold private opinions, formed from their own and occasionally mistaken observations, cooperation under most of these norms becomes fragile\cite{takahashi2006,uchida2010,hilbe2018,perret2021,fujimoto2023}. The difficulty is that a reputation is both the outcome of past assessments and an input to future ones. Under a norm that approves of justified defection, the same refusal to help is judged justified by an observer who regards the recipient as bad and condemned by one who regards the recipient as good\cite{okada2020dilemmas}. Disagreement about the recipient thus turns into disagreement about the donor, which then affects how others act towards the donor and how those actions are judged. An error in judging one individual can thereby propagate to judgements of others\cite{hilbe2018,uchida2013,radzvilavicius2019,fujimoto2022}. To support cooperation, private reputations must therefore meet two requirements\cite{michelmata2024}: they must discriminate between those who cooperate and those who do not, and different observers must largely agree about who is who. Indeed, whether cooperation is stable depends on how strongly the opinions of different observers are correlated\cite{murase2024synchronization}. Agreement alone, however, is not enough: observers who agree that nearly everyone is bad hold reputations that are consistent but uninformative.

% -- 3. Behavioural assessment versus reputation updating
How can private reputations meet both requirements? It helps to distinguish two steps in forming a reputation. Behavioural assessment asks how an action should be evaluated in its social context, which is the question of how to define goodness\cite{ohtsuki2004}. Reputation updating asks how the new assessment should change the observer's existing opinion of the actor. In the standard model, in which reputations are either good or bad, each new assessment simply replaces the previous opinion, so the two steps collapse into one. Many proposed remedies act on assessment. They align the information that observers use, through public institutions\cite{radzvilavicius2021}, gossip\cite{kawakatsu2024}, empathetic evaluation from the donor's perspective\cite{radzvilavicius2019}, synchronized opinions\cite{murase2024synchronization} or the joint assessment of donors and recipients\cite{tham2026}, or they make norms more lenient or probabilistic\cite{okada2020dilemmas,schmid2021,okada2024,murase2023}. Other approaches also change how opinions are updated. Some norms let the assessment depend on the actor's current standing\cite{ohtsuki2004,santos2018} or leave an opinion unchanged when the recipient is in bad standing\cite{sasaki2017,okada2018}. With three reputation levels, an assessment can move a reputation one level at a time instead of switching it outright\cite{tanabe2013,murase2022,yamamoto2025}. Quantitative scores\cite{schmid2023}, records of past behaviour\cite{clark2020} and the aggregation of several observations\cite{michelmata2024} let an opinion reflect more than the latest assessment. The distinction matters because the damage done by misjudgements depends on two factors: how often they occur and how strongly each one moves an established reputation. Assessment mainly governs the first factor, and updating the second. When a single assessment can overturn a reputation, every misjudgement takes full effect and feeds into later assessments; when updating is gradual, isolated misjudgements can be absorbed before they spread. The same norm, applied with the same errors, can therefore produce very different reputations. Reputation updating is thus a consequential dimension of indirect reciprocity that can be studied with the norm held fixed. Alongside the classical question of how to define goodness, it raises a second: how should judgements of others be revised as new evidence arrives?

% -- 4. Asymmetric updating: individual judgement, collective consequences, gap
Once updating is treated as a process in its own right, losing and regaining good standing need not follow the same rule. How strongly people revise an impression depends on what they already believe about the person\cite{mendesiedlecki2013neural,mendesiedlecki2013diagnostic}. Negative information about moral character typically weighs more than positive information, partly because immoral acts are seen as more diagnostic of a person's disposition\cite{reeder1979,skowronski1989,mendesiedlecki2013diagnostic}. Trust, similarly, is said to be easier to destroy than to build\cite{slovic1993}. Yet asymmetries do not always favour condemnation: beliefs about people who have behaved badly are more uncertain, and more readily revised, than beliefs about people who have behaved well\cite{siegel2018}. Models of indirect reciprocity have also considered asymmetries. Upward and downward moves between reputation levels can differ when reputations are public\cite{tanabe2013}, judgement biases that weigh good and bad acts unequally can evolve and increase cooperation\cite{rankin2009}, and pessimistic priors or a bias towards previous opinions can favour cooperation\cite{morsky2024,pandula2024}. These findings do not imply any particular updating rule, but they make symmetry an assumption rather than a given. Under private assessment, such asymmetry matters beyond individual judgement. Because reputations guide later actions and assessments, the way each observer withdraws and restores good standing can reshape an entire reputation system. The distortions that updating has to counter also differ between norms\cite{fujimoto2022,fujimoto2023}. Under norms that condemn any refusal to help, such as Scoring and Shunning, discriminators are condemned when they refuse to help individuals they regard as bad, so cooperators accumulate unwarranted negative assessments\cite{okada2020dilemmas}. Under norms that approve of such refusals, such as Simple Standing, defectors can keep or regain good standing whenever an observer happens to regard their recipient as bad. Stern Judging approves of such refusals but condemns helping someone in bad standing, and is therefore exposed to both distortions. Protecting good standing could shield wrongly condemned cooperators, but it would also shelter defectors. Slowing recovery could keep defectors out, but it would also delay the rehabilitation of cooperators misjudged as bad. Which asymmetry, if any, sustains cooperation may therefore depend on the norm. How the separate organization of reputation loss and recovery shapes the discrimination and agreement of private reputations, and through them the evolution of cooperation, remains unclear.

% -- 5. Here we show
Here we develop a model of indirect reciprocity in which reputations have inertia and in which the loss and the recovery of good standing are regulated separately (Fig.~\ref{fig:model}). Individuals judge one another privately in a donation game, using one of four classic norms (Scoring, Shunning, Stern Judging or Simple Standing), and share no information. An established reputation changes only after a run of contrary assessments: a good reputation passes through warning states before it turns bad, and a bad one through recovery states before it turns good. Because the two sequences can differ in length, losing and regaining standing can require different amounts of evidence. Intermediate states can also carry partial consequences, so that incomplete evidence already affects how an individual is treated. Combining analytical calculations with simulations, we find that updating with a norm-specific asymmetry makes private reputations more discriminating and more consistent across observers, whereas symmetric one-step inertia, with a single intermediate state in each direction, does not. The effective asymmetry mirrors the distortions of each norm: it protects good standing under Scoring and Shunning, slows recovery under Simple Standing and acts in both directions under Stern Judging. Updating rules selected on these reputational criteria alone make highly cooperative states stable and reachable from most initial strategy compositions under all four norms. With symmetric one-step inertia, fewer than a quarter of compositions lead to cooperation under Shunning and Stern Judging. More memory is not necessarily better, however: for a fixed number of intermediate states, how they are divided between loss and recovery can decide between no cooperation and high cooperation. Finally, we show that reputational stability and behavioural sensitivity, which conventional updating ties together, can be tuned jointly. Because the evidence needed to overturn an established reputation is set separately from the immediate consequences of each new assessment, reputations can absorb isolated errors while each negative assessment still carries a cost. Our results suggest that cooperation depends not only on how behaviour is judged but also on how judgements are revised. The way individuals update their opinions of others can shape the collective dynamics of social judgement and, with them, the capacity of reputations to sustain cooperation.
